\clearpage
\thispagestyle{empty}
\pdfbookmark[0]{Magnetischer Verein, 1836. Ausgabe und Quellen}{einleitung-1836-sources}
\begin{center}
\Large\textbf{Einleitung zu den Beobachtungen des Magnetischen Vereins von 1836}
\end{center}

Carl Friedrich Gauss schrieb diese Einleitung für den ersten Band der \textit{Resultate aus den Beobachtungen des magnetischen Vereins}, herausgegeben von Gauss und Wilhelm Weber (Göttingen, 1837), S.~3--12. Die KI-Übertragung in modernes Deutsch mit KI-erstelltem Satz folgt der gesammelten Ausgabe in den \textit{Werken}, Band V (Göttingen, 1867), S.~345--351. Die historischen Ortsnamen und das Entfernungsmaß \textit{Meilen} folgen dieser Fassung.

\medskip\noindent\textbf{Historische englische Übersetzung}\par\smallskip
\noindent W.~Francis übersetzte die Einleitung für Richard Taylors \textit{Scientific Memoirs}, Bd.~II (London, 1841), Art.~II, S.~20--25. Professor Lloyd und Major Sabine revidierten die Übersetzung. Die historische englische Fassung steht mit Anmerkungen auch in André Koch Torres Assis (Hrsg.), \textit{Wilhelm Weber's Main Works on Electrodynamics Translated into English}, Bd.~I (Montreal, 2021), Kap.~8, S.~81--84.

\medskip\noindent
\href{https://gdz.sub.uni-goettingen.de/id/PPN236006339}{Gesammelte Ausgabe (1867)}\quad
\href{https://archive.org/details/ResultateausdenBd1Gaus/page/n9/mode/2up}{Erstdruck (1837)}\par\smallskip\noindent
\href{https://archive.org/details/s1id11857360/page/n82/mode/2up}{Francis (1841)}\quad
\href{https://www.ifi.unicamp.br/~assis/Weber-in-English-Vol-1.pdf\#page=84}{Assis (2021)}
\clearpage
